\documentclass{article}

\usepackage[margin=1in]{geometry}
\usepackage{hyperref}
\usepackage{listings}
\usepackage{enumitem}
\usepackage{booktabs}
\usepackage{microtype}

\hypersetup{
  colorlinks=true,
  linkcolor=blue,
  urlcolor=blue,
}

\lstset{
  basicstyle=\ttfamily\small,
  breaklines=true,
  frame=single,
  framesep=4pt,
  xleftmargin=4pt,
  xrightmargin=4pt,
}

\title{Swimplex Developer Guide}
\author{MATH 196 --- SCIAC Swimming Pipeline}
\date{2026}

\begin{document}
\maketitle
\tableofcontents
\newpage

% ---------------------------------------------------------------------------
\section{Overview}

SCIAC does not publish a centralised results database. Hy-Tek Meet Manager
PDFs are the only structured source of meet results; they live on each school's
Sidearm Sports schedule page and disappear or become stale over time. We scrape
them, parse them into a uniform schema, and feed the structured data into an
optimisation model.

\subsection*{What the PDFs look like}

Hy-Tek Meet Manager produces a characteristic two-column layout. Each page has
fixed x-coordinate columns: place, name, age, school, seed time, finals time,
and points. The column boundaries are consistent within a PDF family but vary
slightly between schools and software versions, so we calibrate them empirically
and round floating-point coordinates before comparison.

\subsection*{Overall data flow}

\begin{lstlisting}
scraper.py         --> data/<season>/pdfs/*.pdf

main.py --parse    --> data/<season>/results.json
                       data/<season>/results.csv

process_results.py --> data/<season>/event_rankings_<gender>.json
                       data/<season>/athlete_profiles_<gender>.json

best_performances.py --> data/<season>/best_performances_<gender>.json

ampl_export.py     --> data/<season>/best_performances_<gender>.dat

data_quality.py    --> console report  (+ optional JSON)
\end{lstlisting}

Every script accepts a \texttt{--season YYYY-YY} flag. Gender-specific scripts
also accept \texttt{--gender Men} or \texttt{--gender Women}.

% ---------------------------------------------------------------------------
\section{\texttt{scraper.py}}

All nine SCIAC schools use the Sidearm Sports platform. Schedule pages are
server-rendered HTML, so plain \texttt{requests} plus \texttt{BeautifulSoup}
is sufficient --- no JavaScript execution required.

\subsection*{Indirect PDF links}

PDFs are not directly linked from schedule pages. Sidearm wraps each file in an
HTML viewer page; the actual file sits on a CloudFront CDN. The scraper detects
this situation in two ways:

\begin{enumerate}[noitemsep]
  \item The HTTP \texttt{Content-Type} header returns \texttt{text/html} rather
        than \texttt{application/pdf}.
  \item The response body does not begin with the \texttt{\%PDF} magic bytes.
\end{enumerate}

When either condition triggers, the scraper parses the viewer page to extract
the real CDN URL from the \texttt{s3\_bucket\_path} attribute or the
\texttt{og:url} \texttt{<meta>} tag.

\subsection*{Date filtering}

To restrict downloads to a given season, the scraper walks up the DOM from each
\texttt{<a>} tag that points to a PDF, looking for sibling or ancestor elements
that contain a recognisable date string. Only links whose associated date falls
within the season window (September 1 to April 30) are downloaded.

\subsection*{Safe filenames}

Each downloaded file is prefixed with the team abbreviation (e.g.\
\texttt{CALT\_results.pdf}) to avoid collisions when the same meet PDF appears
on multiple schools' schedule pages. Spaces and special characters are replaced
with underscores.

\subsection*{Season URL rewriting}

Team schedule URLs contain a trailing \texttt{YYYY-YY} segment. The scraper
replaces this segment at runtime so historical seasons can be scraped without
editing \texttt{config.py}.

\subsection*{Rate limiting}

A \texttt{REQUEST\_DELAY} of 1.5 seconds is enforced between HTTP requests.

\subsection*{Command}
\begin{lstlisting}
python3 main.py --scrape --season 2025-26
\end{lstlisting}

% ---------------------------------------------------------------------------
\section{\texttt{parser.py}}

\subsection*{Extraction approach}

We use \texttt{pdfplumber}'s word-level coordinate extraction (not text blocks
or lines). Each word object carries an \texttt{x0}, \texttt{x1}, \texttt{top},
and \texttt{bottom}. We assign words to columns by their \texttt{x0} coordinate
against empirically calibrated boundaries.

\subsection*{Column layout}

\begin{lstlisting}
place  : x0 <  46
name   : x0 in [46,  187)
age    : x0 in [187, 210)
school : x0 in [210, 360)
seed   : x0 in [360, 440)
finals : x0 in [440, 525)
points : x0 >= 525
\end{lstlisting}

\subsection*{Non-obvious parsing issues}

\begin{description}[noitemsep]
  \item[CID ligatures.]
    Some font families encode the letter \texttt{f} as the Unicode private-use
    codepoint \texttt{(cid:976)}. Without correction, ``Butterfly'' appears as
    ``Butter(cid:976)ly''. We apply a CID-to-Unicode map after extraction.

  \item[Floating-point column boundary.]
    \texttt{pdfplumber} returns coordinates as floating-point numbers.
    A word that visually sits in the school column may have
    \texttt{x0 = 209.9999}, which fails the \texttt{>= 210} boundary check.
    All coordinates are rounded to one decimal place before comparison.

  \item[Athlete ID prefix.]
    Some PDFs prepend an integer athlete ID to the name field
    (e.g.\ \texttt{12345 Smith, John}). We strip leading digit sequences with
    a regex before further parsing.

  \item[Gender/year prefix in school field.]
    School fields sometimes read \texttt{W 18 CMS-CA} rather than just
    \texttt{CMS-CA}. We strip the leading gender letter and year code before
    canonicalisation.

  \item[``First Last'' name order.]
    Most PDFs use ``Last, First'' order. At least one PDF
    (\texttt{CLU\_full\_results\_pp\_v\_clu\_v\_westmont.pdf}) uses
    ``First Last'' order with no comma. Name normalisation detects the missing
    comma and inverts the tokens.

  \item[Seed column shift.]
    In some PDFs the seed column shifts left so that a time value lands in
    the school column. The parser detects a time-shaped string in the school
    column and promotes it to the seed field.
\end{description}

\subsection*{Relay parsing --- two-pass approach}

Hy-Tek relay blocks span several lines and must be accumulated before any
swimmer records can be emitted.

\begin{enumerate}[noitemsep]
  \item \textbf{Team row} (place / relay-name / school / time) is stored as a
        \texttt{pending\_relay\_result} with \texttt{relay\_leg=0} and
        \texttt{is\_relay=True}.
  \item \textbf{Swimmer-name row} (\texttt{1) Last, First Age\ \ 2) Last,
        First Age \ldots}) is split on the leg-number prefixes and stored as a
        list of swimmer names.
  \item \textbf{Split row(s)} carry cumulative 50-yard marks and
        parenthesised per-leg deltas. These are accumulated until a
        non-split row arrives.
  \item \textbf{Finalization}: \texttt{\_finalize\_relay} calls
        \texttt{\_extract\_leg\_splits} to derive per-leg times, then appends
        one \texttt{SwimResult} per swimmer (\texttt{relay\_leg = 1\ldots4}).
\end{enumerate}

\subsection*{\texttt{\_extract\_leg\_splits} stride math}

Hy-Tek prints cumulative 50-yard marks; every 50-yard interval except the very
first has a parenthesised delta. The \emph{stride} is
\texttt{leg\_distance // 50}.

\begin{itemize}[noitemsep]
  \item \textbf{200 relay} (stride = 1): leg 1 time = first bare split;
        legs 2--4 = parenthesised deltas at indices 0, 1, 2.
  \item \textbf{400+ relay} (stride $\geq$ 2): all leg-end times come from
        parenthesised deltas; the index for leg $k$ is
        $(stride - 2) + (k - 1) \times stride$.
\end{itemize}

\subsection*{\texttt{relay\_leg} semantics}

\begin{center}
\begin{tabular}{@{}cll@{}}
\toprule
\texttt{relay\_leg} & Meaning & Downstream treatment \\
\midrule
0 & Relay team row & Relay section of event rankings \\
1 & Flat-start leadoff & Treated as an individual swim (no relay label) \\
$\geq$2 & Exchange-start split & Individual section, event suffixed with \texttt{(Relay Split)} \\
\bottomrule
\end{tabular}
\end{center}

% ---------------------------------------------------------------------------
\section{\texttt{main.py}}

\texttt{main.py} is the entry point. It orchestrates the scrape and parse
steps and applies post-parse cleanup.

\subsection*{Season date range}

The season string \texttt{YYYY-YY} is converted to a date window of
September 1 of the first year through April 30 of the second year. This window
is passed to the scraper for date filtering.

\subsection*{SCIAC filter}

After parsing, we discard results whose school does not canonicalise to one of
the nine SCIAC members. Non-SCIAC opponents appear in many dual-meet PDFs.

\subsection*{Deduplication}

The same meet PDF is frequently downloaded from multiple schools' schedule
pages --- SCIAC Championships appears nine times, once per school. We
deduplicate on the tuple
(\texttt{gender}, \texttt{event}, \texttt{name}, \texttt{school},
\texttt{finals}, \texttt{relay\_leg})
before writing \texttt{results.json}.

\subsection*{Commands}
\begin{lstlisting}
python3 main.py --scrape --season 2025-26   # scrape only
python3 main.py --parse  --season 2025-26   # parse only
python3 main.py          --season 2025-26   # both steps
\end{lstlisting}

% ---------------------------------------------------------------------------
\section{\texttt{process\_results.py}}

This script reads \texttt{results.json} and produces per-event rankings and
per-athlete profiles.

\subsection{Name normalisation (\texttt{\_normalize\_name})}

Six steps applied in order:

\begin{enumerate}[noitemsep]
  \item \textbf{Invert ``First Last''} $\to$ ``Last, First'' when no comma is
        present.
  \item \textbf{Strip leading gender prefix}: \texttt{M.Kiss, Jake}
        $\to$ \texttt{Kiss, Jake}.
  \item \textbf{Strip trailing school/year codes}: suffixes matching
        \texttt{WSO}, \texttt{WFR}, \texttt{W18}, \texttt{MSO}, \texttt{MFR},
        \texttt{M21}, and similar patterns.
  \item \textbf{Strip trailing single-letter middle initial} (upper or lower
        case).
  \item \textbf{Strip orphaned trailing comma}: \texttt{Vanluvanee, }
        $\to$ \texttt{Vanluvanee}.
  \item \textbf{Title-case all-caps first names that contain a vowel}:
        \texttt{Zheng, BO} $\to$ \texttt{Zheng, Bo};
        \texttt{Morris, DJ} is left alone because \texttt{DJ} contains no
        vowel.
\end{enumerate}

\subsection{Name alias resolution}

\begin{description}[noitemsep]
  \item[Case-variant resolution.]
    When the same last name appears with different capitalisation
    (e.g.\ \texttt{DeBoom} vs.\ \texttt{Deboom}), the most-frequent form wins
    and all occurrences are rewritten.

  \item[Last-name-only merging.]
    A bare surname (e.g.\ \texttt{Vanluvanee}) is merged into the full form
    \texttt{Vanluvanee, Adam} when exactly one full form exists in the data.

  \item[Manual aliases.]
    A small table of known nickname/typo pairs verified by hand. Entries are
    resolved before the automated passes.
\end{description}

\subsection{School canonicalisation}

22 raw school-name variants map to 9 canonical SCIAC names. CMS has the most
variants because Hy-Tek PDFs frequently prefix it with a gender letter and
graduation year (e.g.\ \texttt{V 19 CMS-CA}).

\subsection{Swimmer vs.\ diver classification}

Results contain no explicit sport flag. We infer the type from each athlete's
event portfolio:

\begin{enumerate}[noitemsep]
  \item Count distinct dive events and distinct non-relay-split swim events.
  \item If only dive events: \texttt{diver}. If only swim events:
        \texttt{swimmer}.
  \item If both: when the minority type accounts for at least one-third of
        the majority count, classify as \texttt{hybrid} and print a warning;
        otherwise assign whichever type has the larger count.
\end{enumerate}

Relay splits (\texttt{relay\_leg} $\geq$ 2) are excluded from the swim count
so that a diver who legs a relay is not misclassified as a swimmer.

This matters because \texttt{ampl\_export.py} filters to \texttt{type="swimmer"}
only --- diving scores are not times and do not fit the optimisation matrix.

\subsection{6-dive vs.\ 11-dive list classification}

Diving scores are not self-documenting; a score of 280 could be from either
the 6-dive or 11-dive format. Classification is per (\texttt{meet\_name},
\texttt{board}) pair:

\begin{enumerate}[noitemsep]
  \item \textbf{Max score $\geq$ 450}: 11-dive. No D3 diver achieves 450 on
        a 6-dive list.
  \item \textbf{Max score $<$ 200}: 6-dive. An 11-dive list rarely falls
        below 200.
  \item \textbf{Ambiguous range 200--450}: fit a 2-component Gaussian Mixture
        Model on all (\texttt{meet}, \texttt{board}) maximum scores,
        optionally augmented with auxiliary-season data via
        \texttt{--aux-seasons}. The higher-mean component corresponds to
        11-dive lists. The GMM posterior probability is stored as a confidence
        score on every dive result. Requires \texttt{scikit-learn}; falls back
        to a midpoint rule if the library is not installed.
\end{enumerate}

\subsection{Outputs}

\begin{itemize}[noitemsep]
  \item \texttt{event\_rankings\_<gender>.json} --- per-event ranked list
        covering individual events, relay splits, and relay team rows.
  \item \texttt{athlete\_profiles\_<gender>.json} --- per-athlete event history
        with every result ranked within the athlete's own season.
\end{itemize}

\subsection{Command}
\begin{lstlisting}
python3 process_results.py --season 2025-26 --gender Men
python3 process_results.py --season 2025-26 --gender Women \
    --aux-seasons 2024-25   # improves GMM fitting for diving
\end{lstlisting}

% ---------------------------------------------------------------------------
\section{\texttt{best\_performances.py}}

Reads \texttt{athlete\_profiles\_<gender>.json} (output of
\texttt{process\_results.py}) and extracts the rank-1 entry for each
(athlete, event) pair --- the athlete's season best.

Output \texttt{best\_performances\_<gender>.json} contains one record per
(athlete, event) with fields \texttt{best}, \texttt{meet}, \texttt{date},
\texttt{place}, and optionally \texttt{relay\_leg} and \texttt{dive\_format}.

\subsection*{Bare integer filtering}

Both this script and \texttt{process\_results.py} reject times that are bare
integers (no colon or decimal point). These are place numbers that the parser
mistakenly picked up from a mis-formatted row --- most commonly the 1650 Yard
Freestyle section of the SCIAC conference PDF. Diving scores are not filtered
because they can legitimately be round numbers.

\subsection*{Command}
\begin{lstlisting}
python3 best_performances.py --season 2025-26 --gender Men
\end{lstlisting}

% ---------------------------------------------------------------------------
\section{\texttt{ampl\_export.py}}

Converts \texttt{best\_performances\_<gender>.json} into AMPL \texttt{.dat}
format for the roster-optimisation model.

\subsection*{AMPL sets and parameters}

\begin{itemize}[noitemsep]
  \item \texttt{IND\_EVENTS} --- 13 individual events.
  \item \texttt{RELAY\_EVENTS} --- 8 relay-split events (leadoff and exchange
        legs treated separately).
  \item \texttt{ATHLETES} --- all swimmers of type \texttt{swimmer};
        divers are excluded.
  \item \texttt{param times} --- a (athletes $\times$ events) matrix of times
        in seconds. Missing entries receive a sentinel value (default
        \texttt{9999.0}).
\end{itemize}

\subsection*{Event identifiers}

AMPL does not accept spaces or slashes in set element names. We use short
AMPL-friendly identifiers such as \texttt{free100} and \texttt{relay\_back50}.

\subsection*{Command}
\begin{lstlisting}
python3 ampl_export.py --season 2025-26 --gender Men
\end{lstlisting}

% ---------------------------------------------------------------------------
\section{\texttt{data\_quality.py}}

Runs ten validation checks on the parsed data and prints a human-readable
report. Pass \texttt{--json} to write machine-readable output instead.

\begin{enumerate}[noitemsep]
  \item \textbf{Suspect times} --- bare integers (no colon or decimal) and
        times whose pace falls outside 15--50 seconds per 50 yards.
  \item \textbf{Suspect dive scores} --- scores outside the range [50, 800].
  \item \textbf{Unrecognised schools} --- school strings that do not map to a
        canonical SCIAC name after canonicalisation.
  \item \textbf{Residual name suffixes} --- W/M+year codes still present in a
        name after normalisation, indicating the normaliser missed a pattern.
  \item \textbf{Missing first name} --- names that reduce to \texttt{Last, }
        (bare surname with trailing comma) after normalisation.
  \item \textbf{Near-duplicate names} --- pairs sharing the same last name with
        a first-name similarity above 0.80, from the same school. Cross-school
        pairs are skipped to avoid false positives.
  \item \textbf{First-Last format names} --- raw names with no comma, indicating
        that inversion may have failed or was not applied.
  \item \textbf{Case-variant duplicates} --- the same name appearing with
        different capitalisation (e.g.\ \texttt{DeBoom} vs.\ \texttt{Deboom}).
  \item \textbf{Event name anomalies} --- unusual relay distances or rows that
        mix Yard and Meter keywords in the same event name.
  \item \textbf{Relay split sanity} --- relay splits whose time exceeds the
        upper bound implied by the pace rule (50 s per 50 yards).
\end{enumerate}

The script also flags duplicate relay team entries --- the same relay team
result appearing more than once after deduplication.

\subsection*{Command}
\begin{lstlisting}
python3 data_quality.py --season 2025-26
python3 data_quality.py --season 2025-26 --json
\end{lstlisting}

\end{document}
